\section{Book commands (quick reference)}
\label{sec:book-commands}

What follows is deliberately plain: every macro and environment you are expected to type, spelled out so you can search, copy, and teach without hunting a GUI. Programmers hide this stuff behind buttons; here it stays in the open because the whole point of the stack is that \emph{you} own the pipeline. Definitions live under \verb|source/latex/|; you do not have to open those files to use the commands. Copy the patterns below into your own \verb|.tex| chapter files and swap the placeholders for real titles and music paths.

\subsection*{Front matter and back matter headings}

\begin{description}
  \item[\texttt{\textbackslash frontMatterChapter\{Title\}}]
  Starts an unnumbered chapter with the thin rule under the title. Reach for it from \verb|main.tex| for the preface, the list of musical examples, and anything else that should look like a front-matter chapter. Add \verb|\addcontentsline{toc}{chapter}{...}| when you also want a table-of-contents line (see \verb|main.tex| for working examples).
  \item[\texttt{\textbackslash backMatterChapter\{Title\}}]
  Matches the front-matter look, but drops in a PDF anchor and a TOC entry automatically. The bibliography hooks this style through \verb|heading=bookbibliography|.
  \item[\texttt{\textbackslash begin\{bookBackMatterText\}} \ldots\ \texttt{\textbackslash end\{bookBackMatterText\}}]
  Slightly smaller type and tighter leading for dense back matter. The glossary print style wraps entries in this environment so the back pages do not feel like a wall of defaults.
\end{description}

\subsection*{Styled lists (prose only)}

\begin{description}
  \item[\texttt{\textbackslash begin\{bookStyledList\}\{Heading\}} \ldots\ \texttt{\textbackslash end\{bookStyledList\}}]
  Numbered list with a bold heading line in house colors. Put the heading in the mandatory argument, then one \verb|\item| per row.
  \item[\texttt{\textbackslash begin\{bookStyledBulletList\}\{Heading\}} \ldots\ \texttt{\textbackslash end\{bookStyledBulletList\}}]
  Same layout as the numbered list, but labels are bullets. Use either environment when you want a callout that still reads like ordinary book prose instead of a slide deck.
\end{description}

\subsection*{Music examples and the list at the front}
\begingroup
\sloppy

\noindent\texttt{\textbackslash listofmusicalexamples}\\
Drops in the unnumbered ``List of Musical Examples'' chapter and fills it from the running list. Call it once from \verb|main.tex|; the template already does.

\noindent\texttt{\textbackslash beginBookExample\{Short title\}}\\
\texttt{\textbackslash finishBookExample}\\
Opens \textbf{Example $n$} with your short title, reserves a list-of-examples line at the correct page, and closes the vertical spacing after the body. Put anything between begin and finish: prose, graphics, or one of the music helpers below.

\noindent\texttt{\textbackslash beginBookExampleMusic}\\
\texttt{\textbackslash finishBookExampleMusic}\\
Lives inside \verb|\beginBookExample| when the body is a \verb|lilypond-book| fragment such as \verb|\lilypondfile{...}|. Systems get a fixed gap, and each system's graphics are centered so a wide snippet cannot wander past the type block. Put those wrapper lines directly in the chapter file that discusses the example rather than hiding them in an extra one-purpose \texttt{.tex} include.

\noindent\texttt{\textbackslash betweenLilyPondSystem}\\
Default hook between LilyPond systems with a fixed vertical skip. \verb|lilypond-book| calls it between systems; you rarely touch it unless you are doing advanced layout inside a confined group.

\noindent\texttt{\textbackslash bookInlineMusicUseBaseline}\\
\texttt{\textbackslash bookInlineMusicUseVerticalCenter}\\
Pick alignment before \verb|\beginBookInlineMusic|. Baseline is for inline music at the start of a paragraph. Vertical center is for inline music after the first line of a paragraph.

\noindent\texttt{\textbackslash beginBookInlineMusic}\\
\texttt{\textbackslash finishBookInlineMusic}\\
Wraps inline \verb|lilypond| fragments so staff images line up with text. Close with \verb|\finishBookInlineMusic| on the same logical fragment. As with block examples, the wrapper usually belongs directly in the chapter file around the LilyPond import.

\noindent\texttt{\textbackslash beginStretchHeightMusicExample}\\
\texttt{\{height\}}\\
\texttt{\textbackslash finishStretchHeightMusicExample}\\
Fixed-height block; glue between systems stretches with \verb|\vfill|. Same behavior as \verb|\beginFixedHeightMusicExample| and \verb|\finishFixedHeightMusicExample|. Pass a height such as \verb|0.42\textheight|.

\noindent\texttt{\textbackslash beginMultiPageStretchMusicExample}\\
\texttt{\{system\}\{vskip\}}\\
\texttt{\textbackslash finishMultiPageStretchMusicExample}\\
Multi-page excerpt: when LilyPond reaches the system index you name, the layout fires \verb|\clearpage| and optional top glue \verb|vskip|. Other system breaks use stretch glue like the fixed-height helper.

\noindent\texttt{\textbackslash incompleteContentNotice}\\
Draft box with rules for placeholder chapters. Strip the macro call from \verb|main.tex|, or from the chapter \verb|\input| chain, once the chapter is real material.

\noindent\textbf{Minimal display example.} Nesting order never changes: example wrapper, then music wrapper, then LilyPond.

\begin{quote}
\small\ttfamily\raggedright
\textbackslash beginBookExample\{Scale degrees in C major\}\\
\textbackslash beginBookExampleMusic\\
\textbackslash lilypondfile\{./chapters/.../snippet.ly\}\\
\textbackslash finishBookExampleMusic\\
\textbackslash finishBookExample
\end{quote}
\endgroup

\subsection*{Chord symbols and Roman numerals in prose}

\noindent\texttt{\textbackslash MusicChord\{key=value, ...\}}\\
Single interface for chord symbols in running text. Required keys: \verb|root=| (pitch letter) and \verb|quality=maj|, \verb|min|, \verb|dom|, or \verb|aug|. Optional keys: \verb|acc=| (\verb|flat| / \verb|sharp|), \verb|seven=| (for example \verb|7|), \verb|alters={9/flat,5/flat}|, \verb|bass=|, and \verb|bass acc=| for slash chords. The parser accepts legacy aliases too, but \verb|bass=...| and \verb|bass acc=...| are the maintainable forms.

\noindent\texttt{\textbackslash RomNum}, \texttt{\textbackslash RomNumS}, \texttt{\textbackslash RomViiO}\\
Small helpers for Roman numerals in harmonic analysis text.

\noindent\texttt{source/lily/jazzchords.ily}\\
Optional LilyPond-side helper for custom lead-sheet chord naming. Include it from a snippet only when you want engraved chord names to use the jazz-style root namer and exception table; the default template does not force that styling on every score.

\noindent\textbf{Concrete chord line.} The following is valid inside a paragraph:

\begin{quote}
\ttfamily\raggedright
... resolves to \textbackslash MusicChord\{root=F, quality=maj, seven=7\} then\\
\textbackslash MusicChord\{root=G, quality=min, seven=7\}\\
\textbackslash MusicChord\{root=C, quality=dom, seven=7\}\\
\textbackslash MusicChord\{root=F, quality=maj, seven=7\}.
\end{quote}

\subsection*{Glossary and bibliography hooks}

\begin{description}
  \item[\texttt{\textbackslash gls\{label\}}]
  Prints the glossary entry for \verb|label| and wires the hyperref jump. Entries live in \verb|glossary.tex|.
  \item[\texttt{\textbackslash printglossary[style=bookdefinition]}]
  Called once from \verb|main.tex|; the \verb|bookdefinition| style applies the back-matter text environment and a simple name\,\textemdash\,description layout.
  \item[\texttt{\textbackslash printbibliography[heading=bookbibliography, title=\{Bibliography\}]}]
  Prints the bibliography with the same unnumbered chapter treatment as other back matter.
\end{description}

\subsection*{What stays in the style files}

Anything not named above (internal list-of-examples plumbing, raw \verb|\@| commands, glossary style internals) belongs to maintainers editing \verb|page-style.tex|, \verb|music-examples.tex|, or \verb|music-chords.tex|. If a pattern repeats in every chapter, add one macro there and call it from prose so chapter files stay readable.
